\chapter*{Self-reflection}
\label{chap:self-reflection}

In this concluding chapter, I reflect upon my first year as a PhD researcher in the MS3 group within the Department of Microelectronics at TU Delft.
Moving abroad and leaving family and close friends behind -- with the quiet fear of potentially severing some of those connections -- was never going to be an easy step.
The initial months were dominated by the practical hurdles of finding a place to call home and adjusting to a new environment, but once settled, I was quickly reminded of why I chose to take this path.

Shortly after starting at the department, it became clear that I had joined a group of exceptional people, in terms of scientific capability but also on a personal level.
The academic environment at TU Delft offers remarkable room for personal and professional growth.
Being embedded in a department with world-class expertise, strong international ties, and direct connections to industry provides a steady stream of fresh perspectives and practical relevance to the work we do.

Over the past year, my learning curve has been steep.
Working alongside professors who are leaders in microwave sensing has naturally broadened my technical understanding of radar engineering.
Equally important, however, has been the challenge of developing broader research skills.
Learning how to manage tight project schedules, allocate my time effectively, and take genuine ownership of my research direction has been just as transformative.
I feel very fortunate to have landed in this position, and I do not take the opportunity for granted.
Being part of this community is something I value deeply -- both for the chance to travel and collaborate with researchers abroad, and for the opportunity to pass things forward by supporting younger students here at TU Delft.

Ultimately, what made the transition to the Netherlands manageable -- and allowed me to focus on the positive aspects of this new chapter -- was that I did not have to make the move alone.
Moving together with my partner and our two companion cats meant carrying a genuine piece of home with us.
Having that emotional support made it far easier to overcome the reluctance of leaving the familiar behind and to fully embrace life here.
